% RESEARCH_DRAFT_SHA256: b0143544b2752068a7d92f8b4553d9ae503327d5efd08e21712e2b24dd927ff3
% Layout and source-reference wrappers expanded for the web renderer.
\documentclass[11pt]{article}
\usepackage[margin=25mm]{geometry}
\usepackage{amsthm}
\usepackage{fontspec,unicode-math}
\setmainfont{Latin Modern Roman}
\setmathfont{Latin Modern Math}
\usepackage{amsmath,amssymb,mathtools,xurl,hyperref}
\hypersetup{colorlinks=true,urlcolor=blue,linkcolor=blue}
\setcounter{secnumdepth}{0}
\setlength{\parindent}{0pt}
\setlength{\parskip}{5pt}
\setlength{\emergencystretch}{3em}
\newcommand{\sref}[2]{\hyperlink{source:#1:#2}{[S#2]}}
\newcommand{\cataloguescope}[1]{\paragraph{Source record.}#1}
\newcommand{\floor}[1]{\left\lfloor#1\right\rfloor}
\newtheorem{proposition}{Proposition}
\newtheorem{lemma}[proposition]{Lemma}
\begin{document}
\paragraph{Research manuscript.} OpenAI \texttt{gpt-6-astra}, reasoning \texttt{ultra}; 11 October 2026.
\paragraph{Preserved source body.} The complete catalogue mathematical text and sources follow with layout and source links expanded, including source qualifications and its historical preparation notes. New research begins at the section headed Research attempt.

% UnsolvedMath ID 30000680; source catalogue code OWR-1454-001
\section{Tangency Handling in Telea Image Inpainting}\label{30000680}
{\small UnsolvedMath ID 30000680 · Computer Science}
\subsection{Definitions and mathematical statement}

\paragraph{Shared notation.}Standard number systems are $\mathbb N,\mathbb Z,\mathbb Q,\mathbb R,\mathbb C$. Specialized notation and index ranges follow the cited source and the definitions below.


Image inpainting fills unknown image values from known surrounding data. In a pixel-ordering method the reconstruction depends on the sequence in which unknown pixels are assigned.
\paragraph{Record scope and qualifications.}
The catalogue title is retained as metadata. Named certificates, contracts, bounds or frameworks in that title are not asserted as results of the original source extract.
The tangency configuration, update rule, and mathematical meaning of resolving its singularity require the report; no new inpainting method is introduced.

Source-record limitation: The tangency configuration and mathematical criterion for resolving its singularity are absent.
\paragraph{Statement recorded in the supplied source.}
In Telea's image-inpainting algorithm, does changing the order in which pixels are processed suffice to address the singular case of tangency? \href{https://doi.org/10.4171/owr/2007/03}{[S2]}
\paragraph{Source status.}
The supplied catalogue labels this source record open; that metadata was checked against the supplied snapshot on 2026-10-11. The label is preserved for provenance and is not an independent verification of present open status or mathematical correctness. The catalogue literature assessment is dated 2026-08-21.
\subsection{Short English statement}
Decide whether reordering pixels resolves the tangency singularity in the cited inpainting algorithm.
\subsection{Sources}
\begin{itemize}
\item[S1] Ulam AI and the UnsolvedMath Contributors, UnsolvedMath: A Curated Collection of Open Mathematics Problems, record 30000680 (OWR-1454-001). Catalogue curation and metadata: CC BY 4.0. The source extract is retained; definitions, English summaries and qualifications are editorial additions.
\url{https://huggingface.co/datasets/ulamai/UnsolvedMath}.
\item[S2] Trends in Mathematical Imaging and Surface Processing (2007), Oberwolfach Reports, DOI 10.4171/owr/2007/03, pp. 139--194
\url{https://doi.org/10.4171/owr/2007/03}.
\item Source-specific Oberwolfach Reports licences, including CC BY-SA 4.0 where applicable, are retained; see THIRD\_PARTY\_NOTICES.md.
\end{itemize}
\paragraph{Preparation and review.}
The original source and UnsolvedMath contributors retain attribution for the source extract. Codex (GPT-6; exact serving revision not exposed) assisted with definitions, English summaries, scope notes and formatting. No new human mathematical authorship or independently verified solution is asserted. Permanent source ID 30000680 and code OWR-1454-001 checked against the complete supplied registry and current repository. Distinct upstream records, including repeated questions and broad or damaged extracts, are retained for complete coverage of the source CS category. A repeated record is not a new independent conjecture. Complete mathematical-equivalence and current-literature review has not been performed. Preparation and final review follow contribution prompts C and D, with source gaps disclosed. Structural and rendering checks do not establish a new theorem or open status.
\label{end:30000680}

\section{Research attempt}
\subsection{Literature, source context and scope}
The full original catalogue body and source qualifications are preserved above. The relevant primary workshop passages were read in the cached publisher report. Historical claims and subsequent attributed results are distinguished from the mathematical work below. A restricted example, standard lemma or known theorem is not counted as new progress on the entire original target. Missing assumptions and repeated or rhetorical extracts are not silently turned into different solved problems.

\subsection{A closed-flow obstruction to removing every tangency by smooth ordering}
The primary passage describes an ordering function $T$ and a guidance direction $c$, with a special case where $\nabla T\cdot c=0$. It suggests changing the ordering, including a generalized eikonal equation. We test the strong proposed route of finding a globally smooth single-valued ordering with no tangency anywhere. This is a specific continuum interpretation, not a complete reconstruction of every pixel update or every possible nonsmooth ordering.
\begin{proposition}
On the annulus $\{(x,y):1<x^2+y^2<4\}$ with nonzero guidance field $c(x,y)=(-y,x)$, every continuously differentiable real function $T$ has a point with $\nabla T\cdot c=0$ on each circle centered at the origin in the annulus.
\end{proposition}
\begin{proof}
Fix $1<r<2$ and define the differentiable periodic function $h(\theta)=T(r\cos\theta,r\sin\theta)$. The chain rule gives
\[
h'(\theta)=\nabla T(r\cos\theta,r\sin\theta)\cdot(-r\sin\theta,r\cos\theta)=\nabla T\cdot c.
\]
Because $h(0)=h(2\pi)$, Rolle's theorem supplies $\theta\in(0,2\pi)$ with $h'(\theta)=0$. This proves the claim on every circle, without assuming any boundary values or a particular equation for $T$.
\end{proof}
The same issue appears in a directed cycle of discrete dependencies. A strict processing order satisfying $T(v_{i+1})>T(v_i)$ along every edge of a cycle would imply $T(v_1)>T(v_1)$ after chaining all inequalities. Thus at least one edge must fail the strict forward-order requirement. For a directed four-cycle, an order following three consecutive edges attains three forward edges, so this upper bound is sharp.
\subsection{Verification and remaining gaps}
All 24 strict orders of a directed four-cycle were enumerated; the largest number of forward edges is three. The continuous obstruction proves that globally removing tangency through a smooth scalar ordering is impossible for this auxiliary closed-flow field. It does not show that the particular source image contains such a field, nor that tangency necessarily makes every reconstruction unusable. Nonsmooth cuts, altered guidance, local ordering, multiple passes and different update rules are not excluded. No complete answer for the report's particular tangency configuration or a verified inpainting algorithm is obtained.

\section{References and source attribution}
Primary workshop report: \url{https://doi.org/10.4171/owr/2007/03}. Relevant passages read: 2007/3 pp. 152--155, Telea ordering, coherence transport, tangency and suggested generalized-eikonal reordering.

All retrievals were checked on 11 October 2026. Original question attribution remains in the preserved Sources section. External theorems are attributed and not claimed as new discoveries or fully verified merely by citation.


\paragraph{Integration review.} Codex (GPT-6; exact serving revision not exposed) reviewed the written arguments and expanded layout and source-reference wrappers. Repository metadata and the production web renderer were checked using Python 3.13 and JavaScript with MathJax 3.2.2. This model-assisted review is not independent human proof verification; no PDF compilation or formal proof checking is asserted.

\section{Completion Estimate}
\noindent\textbf{COMPLETION: 0\%}
No substantiated new progress toward the entire original target is claimed. The worked results and failed route are delimited above; a known result, illustrative model, or auxiliary obstruction is not counted as original-target completion.

\section{Final status and attribution}
\textbf{Attempt status: unresolved.} This describes this research attempt, not a new ruling on the source catalogue's status or on external literature claims.
The mathematics and this self-audit were generated by OpenAI \texttt{gpt-6-astra}, reasoning \texttt{ultra}, on 11 October 2026. Codex (GPT-6; exact serving revision not exposed) orchestrated the batch and contributed formatting/repository review. No human mathematical authorship, independent human proof checking, or formal verification is claimed. The historical preparation paragraphs in the reproduced catalogue body concern that body only.

Contribution prompt A was applied to the complete definition and supplied source record, followed by the adversarial checks of prompt B. Local reproduction files are \texttt{batch016/checks.py}, \texttt{computations.json}, and the per-ID audit JSON. They were not public links at preparation time. Numerical checks are finite and do not prove unrestricted results. No LaTeX engine was found locally; compilation was not performed by this research agent.

\end{document}
